% commentaire.tex — f173r, vv. 876–919 (conversion initiale des travaux antérieurs)
\comm{878}{}{\italique{Hélène}, v. 461 :\citComm{Αἴγυπτος; ὦ δύστηνος, οἷ \difference{πέπλευκ᾽ ἄρα.}}}
\comm{881}{}{\italique{Hélène}, v. 465 :\citComm{Μενελέως --- \difference{ποῦ δῆτ᾽ ἂν εἴη; πότερον ἐκτὸς ἢ 'ν δόμοις;}}}
\comm{886}{}{\italique{Hélène}, v. 466 :\citComm{Γραῦς --- τόδ᾽ ἐστὶν αὐτοῦ \difference{μνῆμα, παῖς δ᾽ ἄρχει χθονός}.}}
\comm{894}{ὡς τὰς γυναῖκας ἐπὶ κλοπῇ τοῦ χρυσίου,}{«~Subtiliser les bijoux des femmes~» ; voilà un chef d'accusation nouveau contre Mnésiloque d'autant plus étonnant que rien jusqu'ici ne l'a laissé présager. 
Le fait que Critylla porte un flambeau indique pourtant qu'elle est présente depuis l'arrivée de la procession. Ment-elle pour éviter d'avoir à raconter la trop longue affaire ? il pourrait y avoir une autre explication. Il semble que seules les femmes tenant un rôle actif lors de ces cérémonies pouvaient porter des bijoux (prêtresses, canéphores, \&c.) ; c'est ce qui ressort notamment d'un passage de \italique{Lysistrata} :
\citComm{στρωμάτων δὲ ποικίλων καὶ /
χλανιδίων καὶ ξυστίδων καὶ /
\textbf{χρυσίων}, ὅσ᾽ ἐστί μοι, /
οὐ φθόνος ἔνεστί μοι πᾶσι παρέχειν φέρειν /
τοῖς παισίν, \textbf{ὁπόταν} τε θυγάτηρ τινὶ \textbf{κανηφορῇ}. \\
Et de surtouts brodés, petites chlamides, xystides, \italique{bijoux}, tout ce que je possède, il ne me coûte pas d'offrir à chacun de l'emporter, pour ses enfants et pour sa fille, \italique{si elle devenait canéphore}.}
En fait, le terme χρυσίον, \italique{bijou}, \italique{ornement} semble avoir parfois un sens imagé chez Aristophane et désigner les \italique{appâts} féminins. Ainsi, lorsque Pisthétaire s'exclame à propos de la belle \& court vêtue joueuse de flûte Procné (\italique{Oiseaux}, v. 669) :\citComm{Πισθέταιρος --- ὅσον δ᾽ ἔχει τὸν χρυσόν, ὥσπερ παρθένος,}il est bien évident qu'il fait l'éloge de sa beauté, son galbe, ou que sais-je encore\nc{Hésychius (χ 792) : χρυσίον· τὸ τῶν παιδίων αἰδοῖον. Henderson (\italique{Maculate Muse}, p. 131) y flaire un jeu de mot entre τὸν χρυσόν et τὸν κυσόν.} et non pas de sa verroterie. Ce blason sera d'ailleurs continué quelques vers plus loin par le chœur (vers 676-684). 
Autre exemple : lorsque Dicéopolis fait venir la jeune canéphore en lui prodiguant ce conseil (vers 257-258) :\citComm{πρόβαινε, κἀν τὤχλῳ φυλάττεσθαι σφόδρα\\
μή τις λαθών σου περιτράγῃ τὰ χρυσία.\\
\francais{Avance, et surtout prends garde que dans la foule\\
on ne te grignote tes bijoux.}}Certes, comme on l'a vu, les canéphores portaient des bijoux ; mais on devine aisément avec Aristophane qu'en voyant cette jeune beauté dont Dicéopolis vient de dire : « Heureux celui qui t'épousera et te procréera des belettes », la foule va ressentir quelque pulsion ne relevant pas uniquement de la cleptomanie.
Le verbe περιτρώγειν paraît d'ailleurs bien impropre pour une image de vol à la tire\nc{Si impropre que l'on en serait tenté de comprendre λαβών à la place de λαθών.}.
Quant à κλοπῇ, Mnésiloque nous a déjà fait remarquer au vers 205 que le verbe κλέπτειν pouvait par euphémisme désigner une activité qui n'a de commun avec le vol que de se passer au plus noir de la nuit.\\
Terminons en indiquant que cette accusation de tentative de viol lors d'une réunion exclusivement féminine se retrouve ailleurs ; ces assemblées féminines à huis-clos et que la rumeur publique voyait comme des séances de débauche sexuelle avaient tout pour attiser la curiosité de quelque luxurieux. C'est à une telle tentative d'un certain Critias qu'un fragment d'Hipponax\nc{Fragment 30 de l'édition au fort utile commentaire d'Olivier Masson ($=$ fragment 42 de l'édition Degani). Je reprends la correction proposée par Masson : δούμῳ à la place de δούλῳ, qui s'appuie sur une correction identique chez Hésychius (δ 2256) :\citComm{δοῦμος· ἡ οἰκία, ἢ τὴν ἐπὶ τὸ αὐτὸ συνέλευσιν τῶν γυναικῶν,\\
doumos : la maison. \italique{Item} : une réunion en commun des femmes.}} fait allusion :\citComm{οὔ μοι δικαίως μοιχὸς ἁλῶναι δοκέει\\
Κριτίης ὁ Χῖος ἐν $ \cup - \times - $ δούμῳ\\
Ce n'est pas selon moi par justice\nc{Je lis : οὔ μοι δίκᾳ ὡς μοιχὸς ἁλῶναι δοκέει κτλ.} qu'en tant que débauché fût saisi\\
Critias de Khios dans $[$ … $]$ une assemblée féminine.}}
\comm{895}{βαϋζα τοὐμὸν σῶμα βάλλουσα ψόγῳ.}{Taillardat (§569) avait observé que le tour βάλλειν τινὰ ψόγῳ, « accabler quelqu'un de reproches » est d'un ton tragique fait sur le modèle de βάλλειν τινὰ λίθοις\nc{Βάλλειν κακοῖς (Sophocle, \italique{Ajax} 1244), αἰτίᾳ (\italique{id}. \italique{Trach.} 940), φθόνῳ (Eur. \italique{Élect.} 902). Cf. ἀτιμίῃσιν ἰάλλειν (Homère, \italique{Odyssée} XIII, 142), λόγοις ἰάπτειν (Sophocle, \italique{Ajax} 501).}. (Et réciproquement, les mots sont bien souvent considérés comme des projectiles, cf. Taillardat §502.) 
De même, la périphrase τοὐμὸν σῶμα pour ἐγώ est d'un style élevé et particulièrement prisée d'Euripide\nc{Voir le \LSJ \ à l'entrée σῶμα II.}.
Pour le texte, voir le commentaire critique.}
\comm{896}{κακορροθοῦσά.}{Le vocabulaire est euripidéen -- du véritable Euripide s'entend. Le verbe κακορροθεῖν, substitut expressif de κακολογεῖν, est un terme qu'on retrouve exclusivement dans ses tragédies -- sauf reprise aristophanesque (\italique{Acharniens}, v. 576).}
\comm{898}{εἰ μὴ Κρίτυλλά γ᾽ Ἀντιθέου Γαργηττόθεν·}{Je signale à l'attention d'un éventuel exégète parcourant ces pages que la réplique de Critylla trouve un écho dans un fragment de la comédie perdue d'Aristophane des Δαιταλεῖς, \italique{Convives}, transmise par Pollux (X, 158 $=$ fr. 226(217) Kassel-Austin) :\citComm{εἰ μὴ δικῶν τε γυργάθους ψηφισμάτων τε θωμοὺς\\
< φέροντες >}Le rapport est flagrant : même formule commençant par «~εἰ μὴ~» ; Κρίτ-υλλά, nom formé du suffixe -υλλα\nc{Le suffixe en -υλλα/ιλλα est courant pour les noms de femme : Κύδιλλα, Ἥρυλλα, ou Ξένυλλα qui est apparue plus haut \&c. (voir \italique{La formation des noms en grec ancien} de Chantraine, p. 252.)} et du radical κρίτ- qui fait entendre le \italique{jugement} a le même champ sémantique que les δικῶν ; Γαργηττόθεν et γυργάθους sont très proches phonétiquement ainsi que -τιθέου (de Ἀντιθέου) et τε θωμοὺς.\\
Athénée (I, 4d) évoque ce même fragment d'Aristophane dans son premier livre ; ce qui nous permet d'en préciser le contexte. Dans un chapitre assez décousu\nc{Le premier livre, duquel le passage est tiré, ne nous est connu que par l'\italique{épitomé} des \italique{Deipnosophistes}, ce qui complexifie encore la délimitation et l'attribution des citations... Voir l'article de Benoît Louyest, «~L'épitomé du \italique{Banquet des sophistes} d'Athénée~», in \italique{Rursus} 8 (2012).} consacré aux propos de table, on comprend qu'il est certainement question de sycophantes -- pique-assiettes et autres croque-lardons --, dont la vie est douce, bonne et grasse :\citComm{τὸν βίον, φησίν, εὐσταθεῖς· οὐκ ἐγχειρογάστορες, γυργάθους ψηφισμάτων φέροντες, Ἀριστοφάνης φησίν}Comme l'oreille grecque percevait comme particulièrement proches les sons `γ' et `κ', faut-il entendre la locution grecque :\citComm{οὐκ ... εἰ μὴ τὴν δίκην τις θεῶν κατώρθωσεν,\\
À moins qu'un dieu ne rétablisse la justice,}conservée par Isocrate\nc{Isocrate, \italique{Orationes}, V (Φίλιππος), 152 ; XV (Περὶ ἀντιδόσεως), 105.} et auquel Aristophane fait vraisemblablement allusion en \italique{Cavaliers} 1360.}
\comm{898}{Γαργηττόθεν·}{La scholie :\citComm{ἀπὸ δήμου τῆς Αἰγηίδος φυλῆς, ἔνθα τὸ Εὐρυσθέως ἀπόκειται σῶμα ἐν τῷ Ἀρατηρίῳ. ἐπίρρημα δὲ τὸ Γαργηττόθεν.\\
D'un dème de la tribu Égéide, c'est là que repose le corps d'Eurysthée dans l'Arêtérion. D'où le nom de `Gargêttothen'.}}
\comm{903}{αἰσχύνομαί σε τὰς γνάθους ὑβρισμένη.}{Rappelons que lors de la séance de travestissement, le pauvre et anciennement barbu Mnésiloque s'était fait flamber et raser la barbe. Sur un fameux cratère\nc{Le « cratère de Würzburg », Martin-von-Wagner-Museum der Universität Würzburg, H 5697. Reproduit en annexe ; cf. la note aux vers 754-757.} représentant la scène d'égorgement de la `petite' de Mika, on le voit en porter de très visibles stigmates. Et le scholiaste de commenter ce vers :\citComm{ἐπειδὴ ξυρηθεὶς ἦν,\\
puisqu'il a été rasé.}}
\comm{904–5}{}{\italique{Hélène}, v. 549 : « \difference{ἔκπληξιν ἡμῖν} ἀφασίαν \difference{τε προστίθης} » ; v. 557 : « τίς εἶ; τίν᾽ ὄψιν \difference{σήν}, γύναι, \difference{προσδέρκομαι}; » ; v. 72 : « ὦ θεοί, τίν᾽ \difference{εἶδον} ὄψιν; \difference{ἐχθίστην ὁρῶ} ».}
\comm{906–12}{}{\italique{Hélène}, v. 558 sq. :\citComm{Ἑλένη --- σὺ δ᾽ εἶ τίς; αὑτὸς γὰρ σὲ κἄμ᾽ ἔχει λόγος.\\
\difference{Μενελέως --- οὐπώποτ᾽ εἶδον προσφερέστερον δέμας.\\
Ἑλένη --- ὦ θεοί· θεὸς γὰρ καὶ τὸ γιγνώσκειν φίλους.}\\
< Μενελέως --- Ἑλληνὶς εἶ τις ἢ ἐπιχωρία γυνή; > (Ce vers a été restitué à la tragédie d'après Aristophane.)\\
Ἑλένη --- Ἑλληνίς· ἀλλὰ καὶ τὸ σὸν θέλω μαθεῖν.\\
Μενελέως --- Ἑλένῃ σ᾽ ὁμοίαν δὴ μάλιστ᾽ εἶδον, γύναι.\\
Ἑλένη --- ἐγὼ δὲ Μενέλεῴ γε σέ· \difference{οὐδ᾽ ἔχω τί φῶ}.\\
Μενελέως --- ἔγνως γὰρ ὀρθῶς ἄνδρα δυστυχέστατον.\\
Ἑλένη --- ὦ χρόνιος ἐλθὼν σῆς δάμαρτος \difference{ἐς χέρας}.}}
\comm{910}{ἐγὼ δὲ Μενέλαον ὅσα γ᾽ ἐκ τῶν ἰφύων.}{Pour le texte et le sens du vers (Ménélas reconnu à ses hardes), voir le commentaire critique.}
\comm{912}{ἐλθὼν σῆς δάμαρτος ἐσχάρας / λαβέ.}{Le vers correspondant de l'\italique{Hélène} (le vers 566) finit par : ἐς χέρας ; Hélène reproche à Ménélas d'avoir trop tardé à rejoindre ses \italique{mains} (ses \italique{bras}, dira-t-on en français). Occasion par trop belle pour Aristophane de renouer avec la tradition populaire qui ne voit en Hélène qu'une pandorienne dévergondée, à l'origine de tant de maux : il fait donc prononcer à Mnésiloque `ἐσχάρας'. Le mot ἐσχάρα désigne proprement le \italique{foyer} ou le \italique{réchaud}\nc{\italique{Acharniens}, v. 888 ; \italique{Guêpes}, v. 938 ; \&c.}. Ici comme ailleurs\nc{\italique{Cavaliers} v. 1286 : κυκῶν τὰς ἐσχάρας avec le commentaire du scholiaste : τὰ χείλη τῶν γυναικείων αἰδοίων.}, le terme peut désigner le sexe de la femme. Jean Taillardat\nc{\italique{Les images d'Aristophane}, §117.} rapproche cette métaphore d'un passage de la \italique{Paix}\nc{Vers 891.} où le mot ὀπτάνιον, \italique{cuisine} est pris comme synonyme de κύσθος.}
\comm{913–5}{λαβέ με λαβέ με πόσι, / περίβαλε δὲ χέρας. / φέρε σὲ κύσω. ἄπαγέ μ᾽ ἄπαγ᾽ ἄπαγέ με.}{Les quatre dochmiaques sont tous intégralement résolus (et Lætitia\nc{\italique{The Songs of Aristophanes}, p. 429.} note la «~\anglais{correption }κύσ\lettreBreve{ω}. ἄπαγε \anglais{harsh and unusual in dochmiacs}~») ; signe d'une accélération -- voire une précipitation en l'occurrence -- du rythme.\\
Aristophane nous présente Hélène pour ce qu'elle est : une nymphomane. (Loin donc de la pseudo-réhabilitation d'Euripide, plus risible que crédible.)
Comment Hélène nous fait-elle entendre les besoins tyranniques qui la tiraillent ? par un procédé vieux comme le monde et qu'Aristophane emploie aussi au début des \italique{Cavaliers} (vers 26) : répéter un bout de phrase, en l'accélérant... et faire entendre ainsi un autre motif, bien différent de l'initial. 
La luronne détachait-elle le `πάγε με' de la fin de son exhortation afin de suggérer le verbe πήγνυμι, \italique{ficher} ?\\
La source d'inspiration : \italique{Hélène}, v. 625 sq. :\citComm{Ἑλένη --- ὦ φίλτατ᾽ ἀνδρῶν Μενέλεως, ὁ μὲν χρόνος\\
παλαιός, ἡ δὲ τέρψις ἀρτίως πάρα.\\
ἔλαβον ἀσμένα πόσιν ἐμόν, φίλαι,\\
περί τ᾽ ἐπέτασα χέρα\\
φίλιον ἐν μακρᾷ φλογὶ φαεσφόρῳ.\\
Μενελέως --- κἀγὼ σέ· πολλοὺς δ᾽ ἐν μέσῳ λόγους ἔχων\\
οὐκ οἶδ᾽ ὁποίου πρῶτον ἄρξωμαι τὰ νῦν.\\
Ἑλένη --- γέγηθα, κρατὶ δ᾽ ὀρθίους ἐθείρας\\
ἀνεπτέρωκα καὶ δάκρυ σταλάσσω,\\
περὶ δὲ γυῖα χέρας ἔβαλον, ἡδονάν,\\
ὦ πόσις, ὡς λάβω.}}
